\chapter{Data Leakage Prevention}
\label{ch:leak}

\section{Splitting methodology}
\textbf{RadHAR}: the official test recordings (58) are kept unchanged. The
remaining 142 training recordings are split 80/20 \emph{per class by recording}
(113 train, 29 validation, fixed seed 0, independent of the model seed).
\textbf{MiliPoint (primary)}: subject-disjoint, with subjects 8, 9, 10 for test,
subjects 6, 7 for validation, and the remaining six subjects for training. Windows
never cross the boundary of a single-action run.
\textbf{Comparison protocols} (MiliPoint only, never used for model development):
the chronological 80/10/10 cut of the MiliPoint code, and a random 80/10/10
split of overlapping windows.

\section{Restrictions}
\begin{itemize}[nosep]
\item \textbf{Pre-processing}: the only data-dependent statistic, the intensity
  mean and standard deviation, is computed on training recordings and reused
  unchanged for validation and test. Gating radius and ROI threshold come from the
  EDA and are not learned.
\item \textbf{Augmentation}: stochastic augmentation and robust-noise
  augmentation are applied only in training mode, from a separate random stream.
  Evaluation samples are deterministic (fixed per-index seed).
\item \textbf{Test-set protection}: the checkpoint is chosen by validation
  macro-F1. The test split is evaluated once per run, after training. No
  hyper-parameter, pre-processing threshold or architecture choice was changed in
  response to a test score. The only data-handling change made after the first
  runs, the validation cap, was motivated by validation behaviour (Issue \#17).
\item \textbf{Exported model}: the seed exported for the demo is chosen by
  validation macro-F1, not by test accuracy.
\end{itemize}

\section{Automated audit}
\texttt{scripts/leakage\_audit.py} recomputes every check from the data and run
files. Its output (\texttt{docs/leakage\_audit.md}) is reproduced in
Table~\ref{tab:audit}.

\begin{longtable}{L{4.6cm}L{6.6cm}L{2.8cm}}
\caption{Automated leakage-audit results.}\label{tab:audit}\\
\toprule
Check & Result & Status \\
\midrule\endfirsthead
\toprule Check & Result & Status \\ \midrule\endhead
\csname @@input\endcsname ../paper/generated/audit_rows.tex 
\bottomrule
\end{longtable}

\section{Findings}
\begin{enumerate}[nosep]
\item RadHAR splits share no recording, no identical recording content and no
  identical frame. Our splits are therefore recording-disjoint.
\item \textbf{RadHAR's official test set is not session-independent.} The ROS
  timestamps show that every test recording was made within 10 minutes of a
  training recording, on the same four days (median gap 29\,s), and the dataset
  provides no subject IDs. Subject-independent generalisation therefore
  \emph{cannot be verified} on RadHAR, and RadHAR accuracy should be read as
  ``new recordings, same people and room''.
\item On MiliPoint, the subject-disjoint split shares no subject and no frame
  between splits. This is the protocol that tests generalisation to new people,
  and it is our primary result.
\item The chronological protocol of the MiliPoint code shares subjects (0 and 8)
  between training and test.
\item The random-window protocol places \mpAuditLeakFrames\% of test frames inside
  training windows, and \mpAuditLeakWins\% of test windows overlap a training
  window (by \mpAuditLeakOverlap\% of the window on average). Accuracies obtained
  this way measure memorisation of near-duplicates; Chapter~\ref{ch:results}
  quantifies the inflation.
\item All completed runs satisfy ``saved checkpoint = epoch of maximal validation
  macro-F1'' (\auditSelConsistent{} of \auditSelRuns{} runs).
\end{enumerate}
